\section{Estrutura e formatação do texto}\label{sec:texto}

\subsection{Seções e parágrafos}

Use \lstinline|\section{...}|, \lstinline|\subsection{...}| e \lstinline|\subsubsection{...}| para criar a hierarquia do texto. Uma linha vazia no arquivo-fonte inicia um novo parágrafo. Os comandos \lstinline|\newpage| e \lstinline|\clearpage| forçam uma nova página; o segundo também posiciona antes dela as figuras e tabelas que ainda estiverem pendentes.

\subsection{Ênfase, estilos e caracteres especiais}

O comando \texttt{\textbackslash emph\{ênfase\}} produz \emph{ênfase}; \lstinline|\textbf{negrito}| produz \textbf{negrito}; \texttt{\textbackslash textit\{itálico\}} produz \textit{itálico}; e \texttt{\textbackslash texttt\{monoespaçado\}} produz \texttt{monoespaçado}. Para escrever símbolos reservados, use uma barra invertida: \lstinline|\%|, \lstinline|\&|, \lstinline|\_|, \lstinline|\#|, \lstinline|\{| e \lstinline|\}|.

Aspas tipográficas são escritas como \lstinline|``texto''|, produzindo ``texto''. Uma quebra de linha pontual pode ser inserida com \lstinline|\\|, mas não deve substituir a separação normal entre parágrafos.

\subsection{Listas}

O ambiente \texttt{itemize} cria uma lista com marcadores:

\begin{itemize}
  \item cada item começa com \lstinline|\item|;
  \item listas podem ser aninhadas;
  \item o recuo e os marcadores são aplicados automaticamente.
\end{itemize}

O ambiente \texttt{enumerate} cria uma sequência numerada:

\begin{enumerate}
  \item escreva o conteúdo;
  \item compile o documento;
  \item verifique o PDF.
\end{enumerate}

Para termos e definições, use \texttt{description}:

\begin{description}
  \item[Comando] instrução iniciada por barra invertida.
  \item[Ambiente] bloco delimitado por \lstinline|\begin{...}| e \lstinline|\end{...}|.
\end{description}

\subsection{Notas e citações longas}

Uma nota de rodapé é inserida no ponto desejado com \lstinline|\footnote{Texto da nota.}|\footnote{Este é o resultado produzido pelo comando de nota de rodapé.}. Para uma citação longa, utilize o ambiente \texttt{quoting}:

\begin{quoting}[leftmargin=4cm]
Este bloco exemplifica uma citação recuada. Acrescente a chamada bibliográfica antes do ponto final, conforme as regras adotadas pelo trabalho \parencite[p.~10]{lamport1994}.
\end{quoting}
